\section{Metodología: Agentic Engineering vs. Vibe Coding}

Peter ofreció en la entrevista un contraste de alta difusión: \textbf{Agentic Engineering} frente a \textbf{Vibe Coding}. No se trata de un juego de palabras, sino de una cuestión de disciplina de ingeniería.

\subsection{Por qué "vibe coding es peyorativo"}

Lo describió como "el modo de después de las tres de la mañana": alta velocidad, baja restricción, y al día siguiente hay que dedicar mucho tiempo a limpiar. En cambio, agentic engineering enfatiza las restricciones previas y los cambios verificables.

\begin{importantbox}{Disciplina mínima de Agentic Engineering}
\begin{itemize}
    \item Definir primero el objetivo y los límites, y solo después delegar en el agent;
    \item Priorizar que el agent lea el código y la documentación antes de modificarlos;
    \item Controlar el cambio de modo con palabras de activación (modo de discusión / modo de ejecución);
    \item Revisar los efectos secundarios después de cada ronda de ejecución, en lugar de mirar solo si la funcionalidad "corre".
\end{itemize}
\end{importantbox}

\subsection{Desarrollo impulsado por voz: el intercambio entre eficiencia y ergonomía}

En la entrevista mencionó el uso intensivo de voice input, al grado de sufrir una pérdida temporal de voz. Esto sugiere que, aunque la fricción de la interacción disminuye, surge una nueva carga de trabajo (dictar de manera continua, gestionar el contexto).

\begin{warningbox}{Voice-first no siempre es más eficiente}
La interacción por voz es eficaz en la etapa de ideation, pero en tareas que exigen alta precisión, el texto sigue siendo más fácil de auditar y de rastrear. Se recomienda usar la voz para "expresar la intención" y trasladar las restricciones clave de vuelta al texto.
\end{warningbox}

\subsection{Resumen de la sección}
El núcleo de Agentic Engineering no es "saber usar mejor el modelo", sino \textbf{hacer que el modelo acelere la ingeniería dentro de límites gobernables}. La velocidad en sí misma no es el objetivo; la iteración sostenible sí lo es.

